% =================================================
% Ученический рабочий лист
% =================================================

\worksheettitle
  {Точки через равные промежутки}
  {Считаем промежутки и находим расстояния}

\studentnameblock

\begin{checkpointbox}[Входная разминка]
  Нарисуйте пять точек подряд через одинаковые промежутки. Сколько точек вы
  нарисовали? Сколько промежутков получилось?
  \writinglines{2}
\end{checkpointbox}

\section{Точки и промежутки}

\pointsvsintervalsfigure

\begin{fillbox}[Дополните вывод]
  Между пятью точками находится \answerblank[18mm] промежутка. Вообще, если
  подряд отмечено несколько точек, то промежутков на \answerblank[18mm]
  меньше, чем точек.
\end{fillbox}

\begin{practicebox}[Задание 1. Посчитайте]
  \begin{enumerate}[label=\alph*),itemsep=0.55em]
    \item Между 3-й и 9-й точками: \answerblank[24mm] промежутков.
    \item Между 8-й и 27-й точками: \answerblank[24mm] промежутков.
    \item Между 15-й и 16-й точками: \answerblank[24mm] промежуток.
    \item Между 41-й и 56-й точками: \answerblank[24mm] промежутков.
  \end{enumerate}
\end{practicebox}

\studentcountstrip

\begin{practicebox}[Задание 2. Прочитайте рисунок]
  \begin{enumerate}[label=\alph*),itemsep=0.55em]
    \item Сколько изображено точек? \answerblank[20mm]
    \item Сколько промежутков между 2-й и 7-й точками? \answerblank[20mm]
    \item Если один промежуток равен \(3\) см, каково расстояние между
      2-й и 7-й точками? \answerblank[35mm]
  \end{enumerate}
\end{practicebox}

\section{Находим расстояние}

\begin{fillbox}[Алгоритм]
  Сначала найдите число \answerblank[34mm], затем умножьте его на
  \answerblank[55mm].
\end{fillbox}

\begin{practicebox}[Задание 3. Равноудалённые точки]
  Точки на прямой отмечены через \(7\) мм. Найдите расстояние:
  \begin{enumerate}[label=\alph*),itemsep=0.65em]
    \item между 5-й и 12-й точками;
      \quad вычисление: \answerblank[60mm]; ответ: \answerblank[28mm];
    \item между 8-й и 27-й точками;
      \quad вычисление: \answerblank[60mm]; ответ: \answerblank[28mm];
    \item между 31-й и 44-й точками;
      \quad вычисление: \answerblank[60mm]; ответ: \answerblank[28mm].
  \end{enumerate}
\end{practicebox}

\clearpage
\begin{thinkbox}[Задание 4. Найдите ошибку]
  Ученик рассуждал так: «От 6-й до 14-й точки восемь точек, поэтому при шаге
  \(5\) см расстояние равно \(8\cdot5=40\) см».

  В чём ошибка? Запишите правильное решение.
  \writinglines{3}
\end{thinkbox}

\section{Обратные задачи}

\reverseexamplefigure

\begin{practicebox}[Задание 5. Найдите длину одного промежутка]
  \begin{enumerate}[label=\alph*),itemsep=0.75em]
    \item Расстояние между 3-й и 7-й точками равно \(24\) см.
      \newline Число промежутков: \answerblank[25mm].
      Длина одного промежутка: \answerblank[45mm].
    \item Расстояние между 11-й и 19-й точками равно \(56\) мм.
      \newline Число промежутков: \answerblank[25mm].
      Длина одного промежутка: \answerblank[45mm].
  \end{enumerate}
\end{practicebox}

\unknownpointfigure

\begin{practicebox}[Задание 6. Найдите номер точки]
  Точки расположены через \(4\) см.
  \begin{enumerate}[label=\alph*),itemsep=0.65em]
    \item На рисунке подпишите номер последней точки.
    \item От 12-й точки вправо прошли \(28\) см. Получили точку с номером
      \answerblank[22mm].
    \item От 30-й точки в сторону уменьшения номеров прошли \(36\) см.
      Получили точку с номером \answerblank[22mm].
  \end{enumerate}
\end{practicebox}

\section{Столбы и деревья}

\postrowfigure

\begin{practicebox}[Задание 7. Столбы вдоль дороги]
  Семь столбов поставлены через \(4\) м один от другого.
  \begin{enumerate}[label=\alph*),itemsep=0.55em]
    \item Сколько промежутков между первым и седьмым столбами?
      \answerblank[22mm]
    \item Каково расстояние от первого столба до седьмого?
      \answerblank[45mm]
  \end{enumerate}
\end{practicebox}

\clearpage
\begin{thinkbox}[Задание 8. Обратная задача про фонари]
  От первого до последнего фонаря \(54\) м. Фонари стоят через \(6\) м,
  причём фонарь есть в каждом конце ряда. Сколько всего фонарей?

  Схема и решение:
  \studentworkline
  Ответ: \answerblank[35mm].
\end{thinkbox}

\section{Чередование}

\alternatingtreesfigure

\begin{practicebox}[Задание 9. Сосны и берёзы]
  Деревья посажены в один ряд через одинаковые расстояния: сосна, берёза,
  сосна, берёза и так далее.
  \begin{enumerate}[label=\alph*),itemsep=0.65em]
    \item Какое дерево стоит на 17-м месте? \answerblank[32mm]
    \item Какое дерево стоит на 28-м месте? \answerblank[32mm]
    \item Сколько берёз среди первых 20 деревьев? \answerblank[24mm]
    \item Сколько сосен среди первых 21 дерева? \answerblank[24mm]
  \end{enumerate}
\end{practicebox}

\Needspace{15\baselineskip}
\begin{thinkbox}[Задание 10. Два условия одновременно]
  В ряду чередуются сосны и берёзы, первой стоит сосна. Расстояние между
  соседними деревьями равно \(3\) м. Найдите расстояние между 7-м и 18-м
  деревьями и определите вид каждого из них.
  \writinglines{4}
\end{thinkbox}

\begin{checkpointbox}[Задание 11. Верно или неверно]
  Поставьте знак «+» или «−».
  \begin{enumerate}[label=\alph*),itemsep=0.55em]
    \item Между 1-й и 10-й точками десять промежутков. \answerblank[12mm]
    \item Между 25-й и 32-й точками семь промежутков. \answerblank[12mm]
    \item Девять столбов, включая крайние, образуют восемь промежутков.
      \answerblank[12mm]
    \item В ряду «сосна, берёза, сосна, берёза, ...» на чётных местах стоят
      сосны. \answerblank[12mm]
  \end{enumerate}
\end{checkpointbox}

\begin{reflectionbox}[Итог]
  Чтобы найти расстояние между точками с номерами, сначала нужно
  \answerblank[100mm].
\end{reflectionbox}
